\documentclass[10pt,a4paper]{amsart}
\usepackage[margin=19mm]{geometry}
\usepackage{amsmath,amssymb,mathtools}
\usepackage[scr=rsfs,cal=euler]{mathalfa}
\usepackage{xfrac}
\usepackage[unicode,hidelinks,bookmarks=false]{hyperref}
\newcommand{\ie}{i.e.\ }
\newcommand{\cf}{cf.\ }
\newcommand{\resp}{respectively\ }
\newcommand{\Subsection}{\subsection}
\providecommand{\nobreakdash}{\nobreak\hbox{-}}
\setlength{\emergencystretch}{2em}
\multlinegap0pt
\usepackage[shortlabels]{enumitem}
\usepackage{amssymb}
\usepackage{tikz}
\newtheorem{lemma}{Lemma}
\newtheorem{theorem}{Theorem}
\newcommand{\Enc}{\operatorname{Enc}}
\newcommand{\F}{\mathbb F}
\newcommand{\wt}{\operatorname{wt}}
\title{On systematicity of linear function-correcting codes}
\author{Duy Ho}
\date{Source: arXiv:2608.29389v1 (CC BY 4.0)}
\begin{document}
\maketitle

\noindent\emph{Source and licence.} On systematicity of linear function-correcting codes. Version arXiv:2608.29389v1 (CC BY 4.0), distributed by arXiv under the Creative Commons Attribution 4.0 International licence, http://creativecommons.org/licenses/by/4.0/, as recorded in the arXiv licence metadata for this exact version and shown on the abstract page https://arxiv.org/abs/2608.29389v1. Full licence notice: \texttt{licenses/arxiv-2608.29389-CC-BY-4.0.txt}.\par\smallskip

\noindent\textit{Editorial scope.} Counted unit: the article's theorem thm:systematic-distance-three (source0.tex lines 1267--1283). The article's complete original proof of it is delivered with it (source0.tex lines 1285--1354, one \texttt{proof} environment of 70 lines). No mathematics is written, repaired or re-derived: the statement and the proof are the article's own lines, in the article's own notation, and no retained line is substituted or re-flowed.  Retained uncounted as the source-local context the proof needs: lines 1226--1241.  Nothing was omitted from any retained span, so the package carries no omission record.  No retained line contains a citation, so the wrapper carries no bibliography and no bibliography entry.  No retained line contains a figure environment, an image-inclusion command or drawing code, so no image is delivered and the package declares no asset dependency.  Editorial formatting only, in addition: an AMS \texttt{amsart} preamble that restates the article's own declarations and macros used by the retained lines, namely the environments \texttt{lemma}, \texttt{theorem} on separate counters, together with the standard packages those lines need.  The retained environments print in this document as Lemma 1, Theorem 1 in order of appearance, whereas the article prints the same items with the numbering of its own full document.  The complete native source of the article as distributed in the arXiv e-print for this exact version is bundled unchanged as \texttt{sources/208846/source0.tex}.
\begin{lemma}
\label{lem:distance-three-rows}
Let $\Enc\colon V\to C\subseteq\F_q^{k+r}$ be a linear encoding,
systematic relative to $\mathcal E$, with
$G_{\mathcal E}(\Enc)=[\,I_k\mid P\,]$, and let
$p_1,\dots,p_k\in\F_q^r$ denote the rows of $P$. Let
$f\colon V\to\F_q^\ell$ be a linear function. If
$\delta_{\Enc}(f)\geq3$, then for all
$i,j\in\operatorname{supp}(f)$:
\begin{enumerate}
\item $\wt(p_i)\geq2$. In particular, $\langle p_i\rangle$ is a
non-coordinate point of $\mathrm{PG}(r-1,q)$.
\item If $\langle f_i\rangle\neq\langle f_j\rangle$, then
$\langle p_i\rangle\neq\langle p_j\rangle$.
\end{enumerate}
\end{lemma}

\begin{theorem}
\label{thm:systematic-distance-three}
Let $f\colon V\to\F_q^\ell$ be a surjective linear function. For
every $r\geq1$, the following are equivalent:
\begin{enumerate}
\item There exists a linear encoding
$\Enc\colon V\to C\subseteq\F_q^{k+r}$, systematic relative to
$\mathcal E$, with $\delta_{\Enc}(f)\geq3$.
\item $\nu_{\mathcal E}(f)\leq\Lambda_q(r)$.
\end{enumerate}
Consequently,
\[
r_{f,\mathcal E}^{\mathrm{sys,lin}}(3)
=
\min\{r\geq1:\nu_{\mathcal E}(f)\leq\Lambda_q(r)\}.
\]
\end{theorem}

\begin{proof}
Let $\Enc$ be as in~(1), with
$G_{\mathcal E}(\Enc)=[\,I_k\mid P\,]$, and let
$p_1,\dots,p_k$ denote the rows of $P$. For each
$L\in\Phi_{\mathcal E}(f)$, fix an index
$i_L\in\operatorname{supp}(f)$ with $\langle f_{i_L}\rangle=L$. 
Define
\[
\psi\colon\Phi_{\mathcal E}(f)\longrightarrow\mathrm{PG}(r-1,q),
\qquad
\psi(L):=\langle p_{i_L}\rangle.
\]
By Lemma~\ref{lem:distance-three-rows}, $\psi$ is an injective map into the set of non-coordinate points of $\mathrm{PG}(r-1,q)$. Hence
$\nu_{\mathcal E}(f)\leq\Lambda_q(r)$.

Conversely, assume that 
$\nu_{\mathcal E}(f)\leq\Lambda_q(r)$, and choose an injection
\[
\gamma\colon\Phi_{\mathcal E}(f)\longrightarrow
\{\text{non-coordinate points of }\mathrm{PG}(r-1,q)\}.
\]

For each point $L\in\Phi_{\mathcal E}(f)$, we fix a vector $h_L\in\F_q^r$
spanning $\gamma(L)$ and an index $i_L\in\operatorname{supp}(f)$
with $\langle f_{i_L}\rangle=L$.

Let $P\in\F_q^{k\times r}$ be the matrix with rows
$p_i, 1 \le i\le k$, defined as follows.
For $i\notin\operatorname{supp}(f)$, set
$p_i:=0$. 
For each $i\in\operatorname{supp}(f)$, let $L:=\langle f_i\rangle$,
and   write $f_i=\lambda_i f_{i_L}$ with $\lambda_i\in\F_q^\times$.
Set $p_i:=\lambda_ih_L$. 

Let $\Enc$ be the linear encoding with
$G_{\mathcal E}(\Enc)=[\,I_k\mid P\,]$.
This encoding is systematic relative to $\mathcal E$ and has redundancy $r$.
We now verify $\delta_{\Enc}(f)\geq3$ by checking
$\wt([w]_{\mathcal E})+\wt([w]_{\mathcal E}P)\geq3$ for all $w$
with $f(w)\neq0$. It is sufficient to check only those $w$ with $\wt([w]_{\mathcal E}) \le 2$. 

\begin{enumerate}[wide=0pt, leftmargin=*, label=Case~\arabic*:]

\item $w=ce_i$ with $c\neq0$. Here $f(w)=cf_i\neq0$ implies
$i\in\operatorname{supp}(f)$; put $L:=\langle f_i\rangle$. Then
$[w]_{\mathcal E}P=c\lambda_ih_L$ has weight $\wt(h_L)\geq2$, since
$\gamma(L)$ is non-coordinate. So
$\wt(\Enc(w))=\wt([w]_{\mathcal E})+\wt([w]_{\mathcal E}P)\geq3$.

\item  $w=\alpha e_i+\beta e_j$ with $\alpha,\beta\neq0$ and
$f(w)\neq0$, say $i\in\operatorname{supp}(f)$, with
$L=\langle f_i\rangle$. 

If $j\notin\operatorname{supp}(f)$, then
$[w]_{\mathcal E}P=\alpha\lambda_ih_L\neq0$.

If
$j\in\operatorname{supp}(f)$ with $\langle f_j\rangle=L$, then
$f(w)=(\alpha\lambda_i+\beta\lambda_j)f_{i_L}\neq0$, which implies
$\alpha\lambda_i+\beta\lambda_j\neq0$, and
$[w]_{\mathcal E}P=(\alpha\lambda_i+\beta\lambda_j)h_L\neq0$. 

If
$j\in\operatorname{supp}(f)$ with
$\langle f_j\rangle=L'\neq L$, then $h_L$ and $h_{L'}$ span the
distinct points $\gamma(L)\neq\gamma(L')$, so
$[w]_{\mathcal E}P=\alpha\lambda_ih_L+\beta\lambda_j h_{L'}\neq0$.
\end{enumerate} 
In each subcase of Case 2, we have $\wt([w]_{\mathcal E}P) \ge 1$, and so $\wt(\Enc(w)) \ge 3$. Hence $\delta_{\Enc}(f)\geq3$.
\end{proof}

\end{document}
